
Efficient attention mechanisms promise computational savings, but mechanistic claims underlying these savings are rarely validated at the component level. This work introduces a sub-hypothesis verification framework that decomposes efficiency claims into testable components with explicit falsification criteria. Applying this framework to temporal dynamic attention, experiments demonstrate 10.2\% FLOP reduction when reducing temporal steps from T=3 to T=2, with perplexity remaining within 0.9\% of the T=3 baseline. However, the proposed convergence mechanism is not supported: attention entropy increases by 0.98\% across temporal steps (from 4.099 to 4.139), indicating attention becomes more diffuse rather than more focused. Additionally, FLOP reduction remains constant at 5.31\% across sequence lengths 128-1024, providing no sequence-length-dependent scaling advantage. These results illustrate that positive efficiency outcomes can coexist with falsified mechanistic explanations, motivating mechanism-level testing alongside end-to-end benchmarks in efficiency research.

